\cleardoublepage%
\phantomsection%
\addcontentsline{toc}{chapter}{Sammanfattning}
\begin{otherlanguage}{swedish}
\chapter*{Sammanfattning}
\markboth{Sammanfattning}{Sammanfattning}

System för elektronisk stödverksamhet (ESM) registrerar vanligtvis en tidsordnad ström av pulsbeskrivande ord, där pulser från flera sändare är blandade och varje sändare kan byta arbetsmod.
Operatörer behöver två grupperingar ur samma blandade lista: vilka pulser som kommer från samma fysiska sändare, och vilka pulser som delar arbetsmod.
Klassiska kedjor sorterar först pulserna per unik radarsändare och klassificerar sedan moder på de återskapade pulstågen.
I en sådan kedja tar modsteget sorteringens hårda tilldelningar för givna, utan någon uppfattning om hur tillförlitliga de är, så ett uppdelat eller förorenat pulståg försämrar tyst indata till modigenkänningen.
De två uppgifterna bygger dessutom på delvis samma egenskaper hos pulserna, till exempel bärfrekvens, pulsbredd och pulsrepetitionsmönster, vilket tyder på att en gemensam representation, som i fleruppgiftsinlärning, kan fungera bättre för båda än två separat utformade steg.
Denna uppsats behandlar därför de två partitionerna som ett gemensamt klustringsproblem.
En transformerkodare för flera uppgifter föreslås, med en gemensam stam och två uppgiftsspecifika grenar som avbildar ett fönster av pulser på en sändarinbäddning och en modinbäddning.
Grenarna tränas med två övervakade kontrastiva förlustfunktioner som bara skiljer sig i hur positiva par väljs: sändaridentitet respektive modetikett.
Vid inferens tillämpas täthetsbaserad klustring (DBSCAN) oberoende på mod- och sändarinbäddningarna, utan att namn tilldelas från en fast katalog.

Kodaren tränas och utvärderas på en syntetisk, fullständigt märkt samling scenarier.
Varje puls har en sändaridentitet som är lokal för inspelningen och en mod ur en katalog med $19$ typer som delas mellan inspelningar.
Inspelningarna delas upp i fönster om $2000$ pulser, och resultat redovisas på $1169$ icke-överlappande testfönster.
Vi jämför det gemensamma träningsmålet med träning för en enskild uppgift, samt vanlig självuppmärksamhet med roterande kodning av ankomsttid och med en parvis fysisk uppmärksamhetsbias.

Vanlig självuppmärksamhet återskapar redan båda partitionerna på de flesta fönster, med ett medelvärde för justerat Rand-index (ARI) på $0{,}948$ för sändare och $0{,}985$ för moder.
Att koda ankomsttid i uppmärksamheten, antingen med RoPE eller med den föreslagna parvisa fysiska uppmärksamhetsbiasen, höjer var för sig sändar-ARI till $1{,}000$ och tar bort den kvarvarande svansen av misslyckade fönster.
En kombination av de två gör det inte.
Med vanlig uppmärksamhet, utan någon kodning av ankomsttid, presterar en kodare tränad enbart för radiosändarsortering bättre än den gemensamma modellen på sändare (ARI $0{,}999$ mot $0{,}948$).
När tid kodas i uppmärksamheten presterar den gemensamma kodaren något bättre än enkeluppgiftsmodellerna med vanlig uppmärksamhet.
Det gemensamma målet är alltså en avvägning med vanlig uppmärksamhet, inte ett gratis extra huvud.
Modulationstyper som hållits utanför träningen och operativa inspelningar lämnas till framtida arbete.

\vspace{1em}
\noindent\textbf{Nyckelord:} transformer, djup klustring, signalspaning, identifiering av radarsändare, modigenkänning, övervakad kontrastiv inlärning.
\end{otherlanguage}
